\chapter{Spectral Estimation}

Chapter~20 defined the theoretical power spectral density of a random process as the
Fourier transform of its ensemble autocorrelation. In practice, however, you never hold
an infinite ensemble --- you hold one finite recording of $N$ samples. It is tempting to
feed those $N$ samples to a DFT, square the magnitude, and call that the spectrum. That naive picture holds only for one
very special kind of signal: a clean, deterministic, perfectly periodic waveform
that we happened to capture for a whole number of periods. Real signals are not
like that. Your microphone picks up a singer plus room rumble plus thermal hiss.
Your wrist IMU picks up a gesture plus muscle tremor plus the bus you are riding
on. These signals are partly \emph{random}, they go on far longer than any block
you can grab, and the part you care about is buried in the part you do not.

When a signal is random, its DFT is random too. Run the transform on one block,
get one answer; run it on the next block, get a visibly different answer --- same
underlying process, different numbers. So a single DFT is not ``the spectrum'' of
a random signal. It is one noisy \emph{estimate} of an underlying quantity, the
\emph{power spectral density} (PSD), which tells us how the signal's power is
distributed across frequency \emph{on average}. This chapter shows how to turn
noisy DFTs into trustworthy estimates of that average, and what each method
costs.

\begin{focusbox}
The chapter in one sentence: the squared DFT magnitude (the \emph{periodogram})
is an unbiased but hopelessly \emph{noisy} estimate of the power spectrum --- its
scatter does not shrink as you collect more data --- and the fix is to
\emph{average} many short periodograms, trading frequency resolution for a calmer,
more reliable curve. Everything else (windows, Welch, resolution-versus-variance)
is detail on that one trade.
\end{focusbox}

\section{What ``the spectrum'' even means for a random signal}
\label{ch21:psd-meaning}

Take a deterministic finite-energy signal $x(t)$. Its Fourier transform $X(f)$
exists, and $\abs{X(f)}^2$ is the \emph{energy spectral density}: integrate it
over a band and you get the energy in that band (Parseval's theorem,
Chapter~7). But a random signal that hums along forever has
\emph{infinite} energy and \emph{finite power}, so $X(f)$ does not even exist as an
ordinary function. We need a power-based replacement.

\begin{intuitionbox}
Energy is the total; power is energy per unit time. A pure tone playing forever
delivers infinite total energy but a perfectly steady, finite power. For signals
that never end, ``power per unit frequency'' is the right currency, and the
power spectral density $S_x(f)$ is the ledger: $S_x(f)\,\dd f$ is the average power
carried by the sliver of frequencies between $f$ and $f+\dd f$.
\end{intuitionbox}

Formally, take a \emph{wide-sense stationary} random signal --- one whose
statistics do not drift with time, so its mean is constant and its
autocorrelation
\[
r_x[m] \;=\; \E\!\big[x[n]\,x[n+m]\big]
\]
depends only on the lag $m$, not on $n$ (the wide-sense stationarity model from
Chapter~20). The power spectral density is the DTFT of that
autocorrelation:
\[
\boxed{\;S_x(\omega) \;=\; \sum_{m=-\infty}^{\infty} r_x[m]\, e^{-\jj \omega m}.\;}
\]
This is the \emph{Wiener--Khinchin theorem}. It defines the spectrum of a random
signal as the transform of its \emph{autocorrelation} (a stable, deterministic
function) rather than of the signal itself (which is random and ill-defined).

\begin{whybox}
Why autocorrelation? Because $r_x[m]$ is exactly the deterministic skeleton that
survives the randomness. Any single realization of the signal is a wild,
unrepeatable squiggle, but the \emph{average} of $x[n]x[n+m]$ over all
realizations is a fixed number for each lag. Wiener--Khinchin says to set aside
the individual squiggle, keep its correlation structure, and transform that; the
result is a well-defined spectrum. The power at frequency
$\omega$ is governed by how strongly the signal correlates with time-shifted copies
of itself.
\end{whybox}

\begin{connectionbox}
Autocorrelation extends the inner-product geometry of Chapter~3 to random processes:
$r_x[m] = \E[x[n]x[n+m]]$ is a \emph{statistical} inner product between the
signal and a shifted copy of itself, with expectation playing the role of the sum
in $\inner{x}{y}$. And Chapter~20's result that white noise has a flat
spectrum is just Wiener--Khinchin applied to $r_x[m]=\sigma^2\delta[m]$: the
transform of a single spike at lag $0$ is a constant, so white noise spreads its
power equally over all frequencies. That is literally what ``white'' means.
\end{connectionbox}

\subsection{The two difficulties}

We have a clean target, $S_x(\omega)$. Two problems stand in the way of
estimating it.

\begin{enumerate}
\item \textbf{We never see the expectation.} $r_x[m]$ is an average over infinitely
many imaginary repetitions of the experiment. We have \emph{one} finite recording.
We must estimate $r_x[m]$ (and hence $S_x$) from a single chunk of data.
\item \textbf{We have finitely many samples.} The sum in Wiener--Khinchin runs over
all lags $m\in\Z$; we can only form a finite number of them, and the high lags are
estimated from very few sample pairs, so they are themselves garbage.
\end{enumerate}

Estimating $S_x(\omega)$ from a finite, noisy record is \emph{spectral estimation}.
No estimator is perfect; each makes a known trade-off. We start with the most
important estimator, see where it fails, and then fix it.

\section{The periodogram: the obvious estimator}
\label{ch21:periodogram}

The obvious approach is to take your $N$ samples, transform them, square the
magnitude, and normalize. That is the \emph{periodogram}.
\[
\boxed{\;\hat{S}^{\text{per}}_x(\omega_k) \;=\; \frac{1}{N}\,\abs[\Big]{\sum_{n=0}^{N-1} x[n]\,e^{-\jj \omega_k n}}^{2}
\;=\; \frac{1}{N}\,\abs{X[k]}^2.\;}
\]
The $1/N$ converts ``energy of the block'' into ``power'' (energy per sample). The
name comes from Arthur Schuster, who in the 1890s used it to hunt for hidden
periodicities in sunspot and weather records --- he was literally looking for
\emph{periods}, hence ``periodogram''.

\begin{happeningbox}
Algebraically, squaring the DFT magnitude is equivalent to
transforming the signal's \emph{sample} autocorrelation estimated from that block. In
symbols, $\frac1N\abs{X[k]}^2$ is the DFT of
$\hat r_x[m]=\frac1N\sum_n x[n]x[n+m]$. So the
periodogram is Wiener--Khinchin with the true ensemble $r_x[m]$ replaced by
a single-block estimate --- and noisy high-lag terms in $\hat r_x[m]$ drive every
problem that follows.
\end{happeningbox}

\subsection{It is (nearly) unbiased}

On average, the periodogram points at the right answer. For a long enough block,
\[
\E\!\big[\hat{S}^{\text{per}}_x(\omega)\big] \;\approx\; S_x(\omega).
\]
More precisely the expected periodogram is the true PSD \emph{smeared} by the
window's spectrum (we will see exactly how in Section~\ref{ch21:leakage}); as
$N\to\infty$ that smear narrows and the bias vanishes. So if you could average
infinitely many periodograms, you would converge to the truth.

\subsection{Its variance does not go away}

The variance is the problem. For Gaussian data, the periodogram at each frequency behaves
like the true PSD multiplied by an exponential ($\tfrac12\chi^2_2$) random variable with mean $1$ and standard deviation
$\approx 1$:
\[
\hat{S}^{\text{per}}_x(\omega) \;\approx\; S_x(\omega)\cdot \tfrac12\chi^2_2,
\qquad \Var\!\big[\hat{S}^{\text{per}}_x(\omega)\big] \;\approx\; S_x(\omega)^2 .
\]
This variance equals the \emph{square of the quantity we are trying
to measure}, and it \textbf{does not shrink as $N$ grows}.
Collect ten times more data and the periodogram does not get smoother. It just
gets \emph{finer} (more frequency bins), each bin still rattling around its true
value with a standard deviation as big as the value itself.

\begin{trapbox}
``More samples $\Rightarrow$ better spectrum'' is false for the raw periodogram.
More samples buy you more frequency \emph{bins}, not less \emph{scatter} per bin.
A 1-million-point periodogram of white noise is a million-point hairy mess whose
spikes are exactly as tall, relative to the mean, as a 256-point one. This is the
single most common misconception in practical spectral analysis. The estimator is
\emph{not consistent}: it never settles down.
\end{trapbox}

\begin{whybox}
Why does the noise refuse to average out? Because each DFT bin is a sum of $N$
terms, and the central limit theorem makes the real and imaginary parts of $X[k]$
roughly Gaussian. Their squared sum is a chi-squared with $2$ degrees of freedom
--- \emph{regardless of $N$}. Adding more samples adds more bins, each still built
from its own two Gaussian draws. You are not pooling information about a given
frequency; you are spreading a fixed amount of per-frequency information ever more
thinly. To actually reduce the scatter at a frequency you must combine
\emph{several independent looks} at that same frequency, which is what the next
section does.
\end{whybox}

Figure~\ref{ch21:fig-perio-vs-avg} shows the problem and its cure.

\begin{figure}[H]
\centering
\begin{tikzpicture}
\begin{axis}[
  width=12.5cm, height=5.6cm,
  xlabel={normalized frequency $f/f_s$}, ylabel={power (a.u.)},
  xmin=0, xmax=0.5, ymin=0, ymax=2.6,
  axis on top, tick align=outside,
  legend style={font=\small, at={(0.98,0.97)}, anchor=north east},
  title={Raw periodogram (jagged) vs.\ averaged estimate (smooth), flat white noise},
  title style={font=\small},
]
\addplot[sigblue, thin, sharp plot] coordinates {
(0.00,1.42)(0.02,0.21)(0.04,2.30)(0.06,0.55)(0.08,1.80)(0.10,0.12)
(0.12,2.05)(0.14,0.78)(0.16,0.33)(0.18,1.95)(0.20,0.62)(0.22,1.50)
(0.24,0.25)(0.26,2.20)(0.28,0.48)(0.30,1.30)(0.32,0.18)(0.34,1.88)
(0.36,0.70)(0.38,2.40)(0.40,0.40)(0.42,1.10)(0.44,0.30)(0.46,1.70)
(0.48,0.58)(0.50,0.95)
};
\addlegendentry{single periodogram}
\addplot[specred, very thick] coordinates {
(0.00,1.05)(0.05,0.97)(0.10,1.02)(0.15,0.99)(0.20,1.01)(0.25,0.98)
(0.30,1.03)(0.35,1.00)(0.40,0.99)(0.45,1.01)(0.50,1.00)
};
\addlegendentry{average of 32 blocks}
\addplot[gray, dashed, thick] coordinates {(0,1)(0.5,1)};
\addlegendentry{true PSD}
\end{axis}
\end{tikzpicture}
\caption{White noise has a flat true PSD (dashed). A single periodogram (thin blue)
swings between near zero and more than double the truth at every bin, and adding
samples will not calm it. Averaging $32$ independent periodograms (thick red)
hugs the true value: the scatter has shrunk by roughly $\sqrt{32}$.}
\label{ch21:fig-perio-vs-avg}
\end{figure}

\section{The cure: average independent looks (Bartlett)}
\label{ch21:bartlett}

This points to the fix. If one periodogram has standard
deviation as big as the signal, average $K$ \emph{independent} periodograms and the
standard deviation drops by $\sqrt{K}$ --- the standard $1/\sqrt{K}$ averaging law
(Chapter~20). Bartlett (1948) made this concrete: chop the record into
$K$ non-overlapping segments of length $L$ (so $N = KL$), periodogram each, and
average them frequency by frequency.
\[
\boxed{\;\hat{S}^{\text{Bart}}_x(\omega_k)
\;=\; \frac{1}{K}\sum_{i=0}^{K-1} \frac{1}{L}\,\abs[\Big]{\sum_{n=0}^{L-1} x_i[n]\,e^{-\jj\omega_k n}}^2,\;}
\]
where $x_i[n]=x[iL+n]$ is the $i$th segment. The variance falls like
\[
\Var\!\big[\hat{S}^{\text{Bart}}_x(\omega)\big] \;\approx\; \frac{1}{K}\,S_x(\omega)^2 .
\]

\begin{intuitionbox}
Think of each segment as one photograph of a fidgety subject. One photo is blurry
with motion; thirty photos averaged reveal the steady pose underneath. Bartlett's
method takes $K$ snapshots of the spectrum and averages them. The fidget (variance)
washes out; the pose (the true PSD) emerges. The cost shows up in the segment
length $L$.
\end{intuitionbox}

The trade-off is as follows. You have a fixed budget of $N$
samples. You must split it as $N=KL$.

\begin{itemize}
\item \textbf{More segments $K$} (each shorter) $\Rightarrow$ smoother estimate,
lower variance.
\item \textbf{Longer segments $L$} (fewer of them) $\Rightarrow$ finer frequency
resolution, because a length-$L$ DFT can distinguish frequencies only $\approx
f_s/L$ apart.
\end{itemize}

You cannot have both. With $N$ fixed, pushing $K$ up pushes $L$ down and vice
versa. This is the \emph{resolution--variance trade-off}, and it is the central
tension of all spectral estimation. We will give it its own section
(Section~\ref{ch21:resolution-variance}) once we have the last ingredient,
windowing.

\begin{example}[Reading the trade-off off real numbers]
\label{ch21:ex-budget}
Suppose $f_s = 1000$ Hz and you collect $N = 8192$ samples (about $8.2$ s).
Two extreme plans:
\begin{itemize}
\item One block, $L=8192$: resolution $\Delta f = f_s/L \approx 0.12$ Hz ---
exquisite --- but $K=1$, so the curve is the full hairy periodogram.
\item $K=64$ blocks of $L=128$: resolution $\Delta f = f_s/L \approx 7.8$ Hz ---
coarse --- but variance down by $64\times$, standard deviation down by $8\times$.
A clean, readable curve that cannot separate two tones closer than $\sim 8$ Hz.
\end{itemize}
A sensible middle, $K=16$, $L=512$: $\Delta f\approx 2$ Hz and a $4\times$ drop in
standard deviation. Which to pick depends entirely on the question you are asking
of the data.
\end{example}

\section{Windowing and leakage, revisited for estimation}
\label{ch21:leakage}

We met spectral leakage in Chapter~8: chopping a signal to a finite block
multiplies it by a rectangular window, and multiplication in time is convolution in
frequency, so every true spectral line gets convolved with the window's transform
--- a sinc-like shape with a tall main lobe and a long train of side lobes. For
\emph{estimation} this matters in two distinct ways, and keeping them separate is
the key to using windows well.

\begin{figure}[H]
\centering
\begin{tikzpicture}
\begin{axis}[
  width=12.5cm, height=5.4cm,
  xlabel={frequency offset from a true tone (bins)}, ylabel={response (dB)},
  xmin=-8, xmax=8, ymin=-70, ymax=3,
  tick align=outside, axis on top,
  legend style={font=\small, at={(0.98,0.97)}, anchor=north east},
  title={Window spectra: rectangular leaks far; Hann trades width for low side lobes},
  title style={font=\small},
]
\addplot[sigblue, thick] coordinates {
(-8,-40)(-7,-30)(-6,-38)(-5,-27)(-4,-35)(-3,-24)(-2.5,-13)(-2,-30)
(-1.5,-7)(-1,-2)(-0.5,-0.2)(0,0)(0.5,-0.2)(1,-2)(1.5,-7)(2,-30)
(2.5,-13)(3,-24)(4,-35)(5,-27)(6,-38)(7,-30)(8,-40)
};
\addlegendentry{rectangular}
\addplot[specred, thick] coordinates {
(-8,-62)(-7,-58)(-6,-55)(-5,-50)(-4,-47)(-3.5,-31)(-3,-40)(-2.5,-31)
(-2,-25)(-1.5,-12)(-1,-4)(-0.5,-0.8)(0,0)(0.5,-0.8)(1,-4)(1.5,-12)
(2,-25)(2.5,-31)(3,-40)(3.5,-31)(4,-47)(5,-50)(6,-55)(7,-58)(8,-62)
};
\addlegendentry{Hann}
\end{axis}
\end{tikzpicture}
\caption{The fundamental window trade. The rectangular window has the narrowest
main lobe (best resolution) but side lobes only $-13$ dB down that fall off slowly
--- a weak tone can be drowned by a strong neighbor's leakage. The Hann window
roughly doubles the main-lobe width (worse resolution) but pushes side lobes to
$-31$ dB and beyond, so it hides weak components far less. There is no window that
wins on both.}
\label{ch21:fig-windows}
\end{figure}

\subsection{Leakage effect one: bias and the smeared truth}

In expectation, the periodogram equals the true PSD convolved with
the window's \emph{power} spectrum $\abs{W(\omega)}^2$:
\[
\E\!\big[\hat{S}^{\text{per}}_x(\omega)\big] \;=\; \frac{1}{N}\,\big(S_x \conv \abs{W}^2\big)(\omega).
\]
Two consequences. The \emph{main lobe} blurs nearby detail together (a resolution
loss --- you cannot resolve two tones inside one main-lobe width). The \emph{side
lobes} let power from strong components splatter into far-away bins where they do
not belong (a dynamic-range loss --- a quiet tone next to a loud one disappears
under the loud one's skirts). Choosing a window means choosing how to balance these
two losses.

\subsection{Leakage effect two: tapering for the average}

When we average overlapping segments (next section), the window does a second job:
it tapers each segment smoothly to zero at its ends. This matters because a
rectangular cut leaves jagged discontinuities at the segment edges, which dump
energy into the side lobes \emph{and} make overlapping segments share a lot of the
same edge artifacts. A smooth taper (Hann, Hamming, Blackman) softens the edges so
overlapping segments are closer to statistically independent, which is what makes
the averaging actually deliver its variance reduction.

\begin{focusbox}
Two windows you will reach for ninety percent of the time:
\begin{itemize}
\item \textbf{Hann} ($w[n]=0.5-0.5\cos(2\pi n/(L-1))$): the safe default. Good
side-lobe suppression ($-31$ dB), fast roll-off, works well with $50\%$
overlap. Use it unless you have a reason not to.
\item \textbf{Rectangular} (no window): only when you have captured exactly an
integer number of periods of a periodic signal, or when you genuinely need the
sharpest possible resolution and the dynamic range is small.
\end{itemize}
Reach for \textbf{Blackman--Harris} (side lobes below $-90$ dB) when a tiny tone
sits next to a huge one, and accept its fat main lobe.
\end{focusbox}

\begin{connectionbox}
This is the same convolution theorem from Chapters~11 and~14:
multiply in time $\Rightarrow$ convolve in frequency. Windowing applies that
theorem deliberately: the window's transform \emph{is} the smearing kernel, and
window design means shaping a kernel with a narrow
main lobe and small side lobes --- a goal the uncertainty principle forbids you
from achieving perfectly, exactly as it forbids a signal from being both
time-limited and band-limited.
\end{connectionbox}

\section{Welch's method: overlap, window, average}
\label{ch21:welch}

Welch (1967) combined both ideas --- windowing and averaging --- and added one more
trick: \emph{overlapping} the segments. His method is the default spectral
estimator in essentially every toolbox (\texttt{scipy.signal.welch},
MATLAB's \texttt{pwelch}), so it deserves a careful walk-through.

\begin{walkbox}
Why overlap the segments at all? Bartlett's non-overlapping segments waste the
data that the window tapers to zero near each segment's edges --- those samples
barely contribute. By sliding the window forward by only half a segment ($50\%$
overlap), the samples that were down-weighted as one segment's edge become the
well-weighted center of the next. You get \emph{more} segments out of the same $N$
samples, hence more variance reduction, almost for free. The catch: overlapping
segments are not fully independent, so $K$ overlapping segments reduce variance by
somewhat less than a factor of $K$ --- but still more than the smaller number of
non-overlapping segments would.
\end{walkbox}

The recipe, in full (Figure~\ref{ch21:fig-welch-pipeline} shows the flow):

\begin{enumerate}
\item Choose a segment length $L$ (sets resolution $\approx f_s/L$) and a window
$w[n]$ of length $L$ (Hann is the usual pick).
\item Slide the window across the data in steps of $L/2$ (50\% overlap), producing
$K$ overlapping segments $x_i[n] = x[iL/2 + n]$.
\item Window and transform each: $X_i[k] = \sum_{n=0}^{L-1} w[n]\,x_i[n]\,e^{-\jj
2\pi kn/L}$.
\item Form each segment's \emph{modified periodogram} and average:
\[
\boxed{\;\hat{S}^{\text{Welch}}_x(\omega_k) \;=\; \frac{1}{K}\sum_{i=0}^{K-1}
\frac{1}{L\,U}\,\abs{X_i[k]}^2, \qquad U=\frac{1}{L}\sum_{n=0}^{L-1} w[n]^2.\;}
\]
\end{enumerate}

\begin{figure}[H]
\centering
\begin{tikzpicture}[>=Stealth, font=\small, scale=0.88, every node/.style={transform shape}]
  % Full record bar
  \draw[fill=gray!12, draw=gray, rounded corners=2pt] (0, 4.2) rectangle (8.0, 4.7);
  \node[font=\footnotesize] at (4.0, 4.45) {Full data record: $N$ samples};
  \draw[<->, gray] (0, 4.9) -- (8.0, 4.9) node[midway, above, font=\scriptsize] {$N$};

  % Overlap bracket above segments
  \draw[<->, thick, specred] (1.6, 3.7) -- (3.2, 3.7) node[midway, above, font=\scriptsize] {$50\%$ overlap ($L/2$)};

  % Segments with 50% overlap Hann windows (L = 3.2, hop = 1.6)
  % Segment 1: [0, 3.2]
  \begin{scope}[shift={(0, 2.5)}]
    \draw[fill=sigblue!15, draw=sigblue, thick, domain=0:3.2, samples=30]
      plot (\x, {0.9*sin(deg(pi*\x/3.2))^2}) -- cycle;
    \node[sigblue, font=\tiny\bfseries] at (1.6, 0.35) {Seg 1 ($L$)};
  \end{scope}

  % Segment 2: [1.6, 4.8]
  \begin{scope}[shift={(1.6, 1.5)}]
    \draw[fill=targetgreen!15, draw=targetgreen, thick, domain=0:3.2, samples=30]
      plot (\x, {0.9*sin(deg(pi*\x/3.2))^2}) -- cycle;
    \node[targetgreen, font=\tiny\bfseries] at (1.6, 0.35) {Seg 2 ($L$)};
  \end{scope}

  % Segment 3: [3.2, 6.4]
  \begin{scope}[shift={(3.2, 0.5)}]
    \draw[fill=quotientorange!15, draw=quotientorange, thick, domain=0:3.2, samples=30]
      plot (\x, {0.9*sin(deg(pi*\x/3.2))^2}) -- cycle;
    \node[quotientorange, font=\tiny\bfseries] at (1.6, 0.35) {Seg 3 ($L$)};
  \end{scope}

  % Segment 4: [4.8, 8.0]
  \begin{scope}[shift={(4.8, -0.5)}]
    \draw[fill=structurepurple!15, draw=structurepurple, thick, domain=0:3.2, samples=30]
      plot (\x, {0.9*sin(deg(pi*\x/3.2))^2}) -- cycle;
    \node[structurepurple, font=\tiny\bfseries] at (1.6, 0.35) {Seg 4 ($L$)};
  \end{scope}

  % FFT Periodogram blocks
  \node[draw, rounded corners=2pt, fill=blue!5, font=\scriptsize, align=center] (fft1) at (9.4, 2.9) {FFT $\to |X_1|^2$};
  \node[draw, rounded corners=2pt, fill=green!5, font=\scriptsize, align=center] (fft2) at (9.4, 1.9) {FFT $\to |X_2|^2$};
  \node[draw, rounded corners=2pt, fill=orange!5, font=\scriptsize, align=center] (fft3) at (9.4, 0.9) {FFT $\to |X_3|^2$};
  \node[draw, rounded corners=2pt, fill=purple!5, font=\scriptsize, align=center] (fft4) at (9.4, -0.1) {FFT $\to |X_4|^2$};

  \draw[->, thick, sigblue] (3.25, 2.9) -- (fft1.west);
  \draw[->, thick, targetgreen] (4.85, 1.9) -- (fft2.west);
  \draw[->, thick, quotientorange] (6.45, 0.9) -- (fft3.west);
  \draw[->, thick, structurepurple] (8.05, -0.1) -- (fft4.west);

  % Accumulator / Averager block on right
  \node[draw, rounded corners=2pt, fill=blue!8, minimum width=2.4cm, minimum height=3.2cm, align=center, font=\small] (avg) at (12.4, 1.4) {Average\\[4pt]$\frac{1}{K}\sum_{i=1}^K \,\frac{|X_i|^2}{L\,U}$};

  \draw[->, thick] (fft1.east) -- (avg.155);
  \draw[->, thick] (fft2.east) -- (avg.170);
  \draw[->, thick] (fft3.east) -- (avg.190);
  \draw[->, thick] (fft4.east) -- (avg.205);

  \draw[->, very thick] (avg.east) -- ++(1.0, 0) node[right, font=\small] {$\hat{S}^{\text{Welch}}_x(f)$};
\end{tikzpicture}
\caption{Welch's method as a block diagram: the total record of $N$ samples is
divided into overlapping segments of length $L$ (typically $50\%$ overlap). Each segment
is tapered by a window $w[n]$ (such as Hann) to suppress leakage, transformed via FFT,
normalized by the window energy factor $LU$, and averaged across all $K$ segments. The
overlap ensures that samples near segment edges, attenuated by one window, are captured
at the center of the adjacent window.}
\label{ch21:fig-welch-pipeline}
\end{figure}

The factor $U$ is the window's mean-square value, and dividing by it restores the
correct \emph{power} calibration --- without it, a tapering window would make every
estimate read low because it throws away signal energy at the segment edges. Getting
this normalization right is the difference between a PSD you can read absolute power
off of and one that is merely the right shape.

\begin{trapbox}
Two normalization booby-traps. (1) Forgetting the $1/U$ window correction: your PSD
comes out scaled down (by about $2.7\times$ for a Hann window) and your power
readings are simply wrong. (2) Mixing up the \emph{power-spectral-density}
normalization (divide by $f_s$ too, units of power/Hz, for noise-like signals) with
the \emph{amplitude/power-spectrum} normalization (units of power, for discrete
tones). For a sine wave you want the latter so the peak reads the tone's actual
power; for broadband noise you want the former so the level is independent of your
DFT length. \texttt{scipy.signal.welch} exposes this as
\texttt{scaling='density'} vs.\ \texttt{'spectrum'} --- choose deliberately.
\end{trapbox}

\subsection{Welch in code}

\begin{lstlisting}[language=Python, caption={Welch's method by hand, then with SciPy. The two agree.}]
import numpy as np

def welch_psd(x, fs, L=256, overlap=0.5, window='hann'):
    x = np.asarray(x, float)
    step = int(L * (1 - overlap))           # hop between segments
    w = np.hanning(L) if window == 'hann' else np.ones(L)
    U = np.mean(w**2)                        # window power, for calibration
    segs = [x[i:i+L] for i in range(0, len(x) - L + 1, step)]
    psd = np.zeros(L // 2 + 1)
    for s in segs:
        X = np.fft.rfft(w * s)
        psd += (np.abs(X)**2) / (L * U)
    psd /= len(segs)
    psd[1:-1] *= 2                           # fold negative-freq power onto positive
    psd /= fs                                # 'density' scaling: power per Hz
    f = np.fft.rfftfreq(L, 1/fs)
    return f, psd

# Demo: 50 Hz tone in white noise, sampled at 1 kHz
fs = 1000
n  = np.arange(8192)
x  = np.sin(2*np.pi*50*n/fs) + 0.5*np.random.randn(len(n))
f, P = welch_psd(x, fs, L=256)

# Compare against the library (should match closely)
from scipy.signal import welch
f2, P2 = welch(x, fs=fs, nperseg=256, window='hann', scaling='density')
print(np.allclose(P, P2, rtol=1e-6))   # True
\end{lstlisting}

\begin{figure}[H]
\centering
\begin{tikzpicture}
\begin{axis}[
  width=12.5cm, height=5.6cm,
  xlabel={frequency (Hz)}, ylabel={PSD (dB / Hz)},
  xmin=0, xmax=200, ymin=-45, ymax=5,
  tick align=outside, axis on top,
  legend style={font=\small, at={(0.98,0.05)}, anchor=south east},
  title={A 50 Hz tone buried in white noise, revealed by Welch averaging},
  title style={font=\small},
]
\addplot[sigblue, thin] coordinates {
(0,-22)(8,-30)(16,-18)(24,-33)(32,-15)(40,-28)(46,-8)(50,0)(54,-9)
(58,-26)(66,-14)(74,-31)(82,-17)(90,-29)(98,-13)(106,-34)(114,-19)
(122,-27)(130,-16)(138,-32)(146,-12)(154,-30)(162,-20)(170,-25)
(178,-15)(186,-33)(194,-21)(200,-28)
};
\addlegendentry{single periodogram}
\addplot[specred, very thick] coordinates {
(0,-21)(10,-20)(20,-20.5)(30,-20)(40,-19)(46,-9)(50,-1)(54,-9)(60,-19.5)
(70,-20)(80,-20.3)(90,-20)(100,-20.2)(110,-20)(120,-20.1)(130,-20)
(140,-20.2)(150,-20)(160,-20.1)(170,-20)(180,-20.2)(190,-20)(200,-20.1)
};
\addlegendentry{Welch (Hann, $K{=}63$)}
\end{axis}
\end{tikzpicture}
\caption{A single periodogram (thin blue) has a noise floor so jagged
that the $50$ Hz tone is only barely the tallest spike. Welch averaging (thick red)
flattens the floor to a clean, readable $-20$ dB and the tone stands $19$ dB
above it. Both curves come from the same data; Welch trades resolution for lower
variance.}
\label{ch21:fig-welch-tone}
\end{figure}

\begin{historybox}
The quest to measure power spectra began in 1898, when British physicist Arthur Schuster coined
``periodogram'' to hunt for suspected periodicities in sunspot records and terrestrial magnetism.
For half a century, scientists plotted raw periodograms and fell into traps, mistaking random fluctuations
for true periodic cycles. The mathematical cause remained obscure until the 1940s and 1950s, when
Maurice Bartlett, John Tukey, and Ralph Blackman proved that the raw periodogram is an inconsistent
estimator whose variance does not decay with sample size.

In 1967, Peter Welch at IBM published the method that now bears his name. The timing
mattered: James Cooley and John Tukey had published the Fast Fourier Transform algorithm just two years
earlier in 1965, slashing the cost of computing multiple transforms. Welch showed how to split data into
overlapping, windowed segments and average their periodograms, turning raw computational speed into
a statistically reliable estimate. Today, calling \texttt{scipy.signal.welch} runs essentially Welch's
1967 recipe.

Welch's method established the default pipeline across oceanography, EEG analysis, and communications.
Yet dividing a finite record into short segments carries an unavoidable physical trade-off: shorter
segments widen the spectral window's main lobe, smearing closely spaced frequency tones into a single
unresolved hump. Where non-parametric averaging hits this fundamental resolution limit, subsequent
chapters explore how multirate filter banks (Chapter~22) and parametric model-based methods extract
pinpoint spectral lines by fitting autoregressive equations directly to the time series.
\end{historybox}

\section{Resolution versus variance in numbers}
\label{ch21:resolution-variance}

We can now put numbers on the trade-off. Fix
your total record length $N$ and sampling rate $f_s$. The two quantities you care
about pull in opposite directions through the segment length $L$.

\textbf{Frequency resolution.} A length-$L$ windowed DFT cannot distinguish two
tones closer than roughly the window's main-lobe width:
\[
\Delta f \;\approx\; C_w \cdot \frac{f_s}{L},
\]
where $C_w$ is a window-dependent constant: $C_w\approx 0.89$ for rectangular,
$\approx 1.5$ for Hann, $\approx 1.8$ for Blackman. Larger $L$ $\Rightarrow$ finer
$\Delta f$. Good resolution wants \emph{long} segments.

\textbf{Variance.} With $50\%$ overlap you get about $K\approx 2N/L$ segments, and
the estimate's variance falls roughly as $1/K$:
\[
\Var\!\big[\hat{S}_x(\omega)\big] \;\approx\; \frac{S_x(\omega)^2}{K} \;\approx\;
\frac{L}{2N}\,S_x(\omega)^2 .
\]
Smaller $L$ $\Rightarrow$ more segments $\Rightarrow$ lower variance. A smooth
estimate wants \emph{short} segments.

\begin{aestheticbox}
The two formulas point in opposite directions. Resolution improves
as $L$ grows; variance improves as $L$ shrinks. With $N$ fixed you slide a single
knob, $L$, between ``sharp but noisy'' and ``smooth but blurry''. There is no
setting that is best for everything --- the right $L$ is dictated by your question.
Are you trying to \emph{separate} two close tones? Push $L$ up. Are you trying to
\emph{measure the level} of a broadband floor? Push $L$ down. Most practical
spectral-analysis choices come down to this trade-off.
\end{aestheticbox}

\begin{figure}[H]
\centering
\begin{tikzpicture}
\begin{axis}[
  width=12cm, height=5.2cm,
  xlabel={segment length $L$ (samples, log scale)}, ylabel={(normalized)},
  xmode=log, log basis x=2,
  xmin=16, xmax=4096, ymin=0, ymax=1.05,
  tick align=outside, axis on top,
  legend style={font=\small, at={(0.5,1.18)}, anchor=south, legend columns=2},
  title={The two curves you are balancing (fixed $N$)},
  title style={font=\small},
]
\addplot[sigblue, very thick, domain=16:4096, samples=40] {1/sqrt(x/16)};
\addlegendentry{frequency blur $\Delta f$ (want small)}
\addplot[specred, very thick, domain=16:4096, samples=40] {sqrt(x/4096)};
\addlegendentry{estimate variance (want small)}
\end{axis}
\end{tikzpicture}
\caption{Two opposing pulls as the segment length $L$ varies with the total record
$N$ held fixed. Frequency blur (blue) shrinks as $L$ grows --- long segments
resolve fine detail. Estimate variance (red) grows with $L$ --- long segments mean
fewer of them to average. The good operating point sits where the curves cross for
\emph{your} priorities, not at either extreme.}
\label{ch21:fig-tradeoff}
\end{figure}

\begin{remark}[The equivalent-degrees-of-freedom view]
\label{ch21:edof}
Statisticians summarize the smoothness of a Welch estimate by its
\emph{equivalent degrees of freedom}, $\nu \approx 2K$ (each segment contributing a
$\chi^2_2$). The fractional standard error of the estimate is $\approx
\sqrt{2/\nu} = 1/\sqrt{K}$. So ``how many segments'' and ``how reliable is each PSD
value'' are the same question: $K=1$ gives a $100\%$ standard error (the raw
periodogram's hopeless scatter); $K=100$ gives $10\%$. This number, not the pretty
curve, is what you cite when someone asks ``is that peak real?''.
\end{remark}

\section{Worked example: choosing parameters end to end}
\label{ch21:worked}

\begin{example}[Separating two close tones in noise]
\label{ch21:ex-twotone}
You record $N=4096$ samples at $f_s=2000$ Hz. You suspect two tones near $300$ Hz
and $312$ Hz, $12$ Hz apart, plus broadband noise. Design a Welch analysis.

\emph{Resolution requirement.} To separate tones $12$ Hz apart with a Hann window
($C_w\approx 1.5$) we need $\Delta f \le 12$ Hz, i.e.
\[
1.5\,\frac{f_s}{L} \le 12 \;\Rightarrow\; L \ge 1.5\cdot\frac{2000}{12} \approx 250.
\]
Round up to a power of two: $L = 512$ (giving $\Delta f \approx 5.9$ Hz, a
comfortable margin).

\emph{How many segments?} With $50\%$ overlap and step $256$:
\[
K \approx \frac{N - L}{L/2} + 1 = \frac{4096 - 512}{256} + 1 = 15.
\]
So $\nu\approx 30$ degrees of freedom and a fractional standard error of
$1/\sqrt{15}\approx 26\%$ --- enough to see the two peaks clearly and trust they
are real, though the noise floor will still look slightly rough.

\emph{If you needed a calmer floor}, drop to $L=256$: now $\Delta f\approx 11.7$ Hz
--- right at the edge of resolving the two tones, risky --- but $K\approx 31$ and
standard error $\approx 18\%$. The tension is explicit: $L=512$ to \emph{see} the
two tones, $L=256$ to \emph{measure the floor}. With only $4096$ samples you cannot
do both superbly; collecting more data is the only real escape, and it buys you
more segments (lower variance) at whatever resolution you have already fixed.
\end{example}

\begin{labbox}
Generate exactly this signal and watch the trade-off:
\begin{lstlisting}[language=Python]
import numpy as np; from scipy.signal import welch
fs, N = 2000, 4096
n = np.arange(N)
x = (np.sin(2*np.pi*300*n/fs) + np.sin(2*np.pi*312*n/fs)
     + 0.7*np.random.randn(N))
for L in (128, 256, 512, 1024):
    f, P = welch(x, fs=fs, nperseg=L, window='hann')
    # zoom the plot to 280..330 Hz and watch the twin peaks merge as L drops
    print(L, "resolves two peaks:", L >= 256)
\end{lstlisting}
At $L=128$ the two tones fuse into one blob; at $L=512$ they split cleanly but the
floor is rougher; at $L=1024$ the split is sharp yet you have only $\approx 7$
segments, so the floor rattles.
\end{labbox}

\section{Spectral analysis of real audio and motion signals}
\label{ch21:real-signals}

\begin{applicationbox}
\textbf{Audio.} A $4$-second guitar note at $f_s=44.1$ kHz is $176{,}400$ samples. Welch with
$L=4096$ and Hann gives $\Delta f\approx 16$ Hz --- fine enough to separate the
fundamental ($\sim 82$ Hz for low E) from its harmonics at $164$, $246$ Hz, etc.
--- and roughly $K\approx 85$ segments for a smooth, trustworthy harmonic profile.
That PSD is the raw material for pitch detection, timbre fingerprinting, and the
audio-feature pipelines of later chapters. For \emph{music}, though, you often want
the \emph{spectrogram} (short-time Fourier transform) instead of a single averaged
PSD, because the spectrum changes over the note's attack, sustain, and decay ---
Welch deliberately averages that time variation away, which is wrong when the
time variation is what you want to study.
\end{applicationbox}

\begin{applicationbox}
\textbf{Motion / IMU.} A wrist IMU sampled at $f_s=100$ Hz captures human motion that lives almost
entirely below $20$ Hz: gait is $1$--$3$ Hz, voluntary arm gestures $0.5$--$5$ Hz,
physiological tremor $4$--$12$ Hz. To resolve a $5$ Hz tremor from a $3$ Hz gesture
($2$ Hz apart) with a Hann window you need $L \ge 1.5\cdot 100/2 = 75$, say
$L=128$ ($\Delta f\approx 0.8$ Hz, $\sim 1.3$ s per segment). A $20$-second
recording then yields $K\approx 30$ segments --- a clean PSD of the motion's
rhythmic content. The accelerometer PSD typically shows a strong low-frequency
hump (intentional motion) and, in some subjects, a distinct tremor peak; the gyro
PSD complements it for rotational components.
\end{applicationbox}

\begin{figure}[H]
\centering
\begin{tikzpicture}
\begin{axis}[
  width=12.5cm, height=5.4cm,
  xlabel={frequency (Hz)}, ylabel={PSD (dB)},
  xmin=0, xmax=25, ymin=-40, ymax=5,
  tick align=outside, axis on top,
  legend style={font=\small, at={(0.98,0.97)}, anchor=north east},
  title={Welch PSD of wrist-accelerometer data: a gesture hump and a tremor peak},
  title style={font=\small},
]
\addplot[sigblue, very thick] coordinates {
(0,-6)(0.5,-3)(1,-1)(1.5,0)(2,-2)(2.5,-6)(3,-12)(3.5,-18)(4,-22)
(4.5,-19)(5,-14)(5.5,-19)(6,-24)(7,-28)(8,-30)(9,-31)(10,-32)
(12,-33)(14,-34)(16,-34)(18,-35)(20,-35)(22,-36)(25,-36)
};
\addlegendentry{accel magnitude PSD}
\node[specred] at (axis cs:1.5,2.2) {\small motion};
\node[specred] at (axis cs:5,-9) {\small tremor};
\end{axis}
\end{tikzpicture}
\caption{A representative Welch PSD of wrist-worn accelerometer data during a slow
arm gesture. A broad low-frequency hump near $1$--$2$ Hz is the intentional motion;
a smaller distinct peak near $5$ Hz is physiological tremor. The averaging is what
makes the tremor peak stand out cleanly from the noise floor instead of being one
of dozens of random spikes.}
\label{ch21:fig-imu-psd}
\end{figure}

\begin{goalbox}
For the wrist instrument controller, the PSD is a \emph{feature extractor}, not a
display. A short Welch (or simpler: the band-power of a single windowed DFT) over a
sliding $\sim 1$ s window gives you a handful of numbers --- power in $0$--$3$ Hz
(deliberate motion), $4$--$12$ Hz (tremor), and higher bands (taps, impacts) ---
that feed the gesture classifier downstream. These band powers are cheap, robust to
the exact gesture timing, and far more stable than raw samples.

Be clear about what they can and cannot do. A wrist IMU's spectrum can tell you
\emph{that} the arm is moving rhythmically and roughly how fast, and it can flag
tremor or taps. It \textbf{cannot} recover individual finger joint angles --- the
spectral signature of finger motion barely reaches the wrist. Continuous finger
pose needs sensors on the fingers (instrumented rings) or muscle sensing (EMG /
force myography), whose own spectra we estimate with these same Welch tools
in Chapter~37. The PSD is a strong coarse-gesture and rhythm
detector; it is not a finger tracker, and any product claim otherwise is false.
\end{goalbox}

\begin{buildbox}
On an Arduino-class board you usually cannot afford a full Welch with overlapping
buffers in RAM. The lightweight field version: keep one circular buffer of $L=128$
or $256$ samples, apply a precomputed Hann window (store it as a constant array),
run a fixed-point or CMSIS-DSP FFT (\texttt{arm\_rfft\_fast\_f32}), square the bins,
and \emph{exponentially average} the resulting periodogram into a running estimate
\[
\hat S \leftarrow \alpha\,\hat S + (1-\alpha)\,\abs{X}^2, \qquad \alpha\approx 0.9.
\]
This is ``Welch on a budget'': the exponential average plays the role of averaging
$K\approx 1/(1-\alpha)\approx 10$ segments, smoothing the estimate without ever
storing more than one segment at a time. Then sum bins into your $3$--$4$ band
powers and ship those to the classifier. A Cortex-M4 at $80$ MHz does a $256$-point
real FFT in well under a millisecond, leaving plenty of headroom at $100$ Hz IMU
rates.
\end{buildbox}

\section{Beyond the periodogram: a brief map}
\label{ch21:beyond}

Welch is the primary tool, but it is not the only option, and it helps to know what
lies past it so you can recognize when you have outgrown it.

\begin{itemize}
\item \textbf{Multitaper (Thomson).} Instead of one window, apply several specially
designed orthogonal windows (discrete prolate spheroidal sequences) to the
\emph{same} segment and average their periodograms. This squeezes more independent
looks out of a \emph{short} record than Welch can --- the best general-purpose
choice when data is precious (short EEG epochs, brief gestures). The cost is more
computation and a fixed resolution--bandwidth product you set up front.
\item \textbf{Parametric / AR methods (Yule--Walker, Burg).} Assume the signal is
the output of an all-pole filter driven by white noise, fit the filter
coefficients, and read the spectrum off the model. These give very sharp peaks from
short records --- excellent when you \emph{know} the signal is a few resonances
(speech formants, a couple of vibration modes) --- but they hallucinate structure
if the model order is wrong, so they demand more faith in your assumptions than the
assumption-free periodogram family.
\item \textbf{The periodogram family stays the safe default.} When you do not have
strong prior knowledge, Welch (or multitaper for short records) makes the fewest
assumptions and fails the most gracefully. Reach for parametric methods only when
you can defend the model.
\end{itemize}

\begin{recoverybox}
\textbf{What this chapter covered:}
\begin{itemize}
  \item Bartlett's and Welch's segment-averaged estimators tame the non-vanishing variance of the raw periodogram, trading frequency resolution for $1/K$ variance reduction.
  \item Segment length $L$ and window shape give explicit quantitative control over the resolution--variance trade-off ($\Delta f = C_w f_s/L$ versus fractional variance $\propto L/N$).
  \item Dividing by the window power factor $U = \frac{1}{L}\sum w[n]^2$ yields physically calibrated power spectral density estimates in $\text{power}/\text{Hz}$.
\end{itemize}

\textbf{Limitations:}
\begin{itemize}
  \item Non-parametric methods cannot escape the classical uncertainty principle: fine frequency resolution demands long segments, which inevitably limits variance reduction for a finite record of $N$ samples.
  \item Windowing always smears: main-lobe width blurs nearby peaks together, while side-lobe leakage can submerge weak signals near strong tones.
  \item Averaging over time inherently assumes wide-sense stationarity. For transient or non-stationary signals (musical note attacks, speech plosives, sudden gesture impacts), Welch averaging destroys the temporal dynamics.
\end{itemize}

\textbf{Forward roadmap:}
\begin{itemize}
  \item \textbf{Chapter 22 (Decimation, Interpolation, and Resampling):} When signals have narrow bandwidths or need to match disparate audio, sensor, and radio clocks, retaining a fixed, high sampling rate wastes compute. Chapter~22 develops the theory of downsampling, upsampling, and fractional rate conversion with anti-aliasing and anti-imaging filters.
  \item \textbf{Chapter 23 (Modulation and I/Q Signals):} Takes spectral analysis into communication systems, frequency shifting, and analytic complex baseband signals.
  \item \textbf{Chapters 37/38 (Feature Extraction and Hand Gesture Classification):} Translates Welch band-power features from forearm EMG and wrist IMU signals into low-latency feature vectors for gesture recognition pipelines.
\end{itemize}
\end{recoverybox}

\section{Exercises}
\label{ch21:exercises}

\begin{exercise}
\textbf{The variance that won't die.} Generate $N=1024$ samples of white Gaussian
noise and plot its periodogram. Now do the same for $N=16384$. Confirm by eye (and
by computing the ratio of standard deviation to mean across bins) that the relative
scatter is the \emph{same} in both, even though the second has $16\times$ more bins.
Explain in one sentence why this is the expected result, not a bug.
\end{exercise}

\begin{exercise}
\textbf{Averaging law.} For white noise, compute Welch estimates with $K=1,4,16,64$
non-overlapping segments (adjust $N$ so $L$ stays fixed). Measure the standard
deviation of the estimate across frequency for each $K$ and verify it falls like
$1/\sqrt{K}$. By how much would you expect $K=256$ to improve over $K=1$?
\end{exercise}

\begin{exercise}
\textbf{Window dynamic range.} Build a signal with a strong tone at $100$ Hz
(amplitude $1$) and a weak tone at $130$ Hz (amplitude $0.01$, i.e.\ $-40$ dB), in
light noise, at $f_s=1000$ Hz. Estimate the PSD with a rectangular window and with a
Blackman--Harris window at the same $L$. Which window lets you see the weak tone,
and why? Relate your answer to side-lobe level.
\end{exercise}

\begin{exercise}
\textbf{Resolution requirement.} You must separate tones at $440$ Hz and $445$ Hz
($f_s=8000$ Hz) using a Hann window ($C_w\approx 1.5$). What is the minimum segment
length $L$? If you have $N=8000$ samples ($1$ s) and use $50\%$ overlap, how many
segments do you get, and what is the resulting fractional standard error?
\end{exercise}

\begin{exercise}
\textbf{Normalization sanity check.} Take a pure $1$ V-amplitude sinusoid (no noise)
and compute its Welch PSD with \texttt{scaling='spectrum'}. The peak should read the
tone's power, $0.5$ V$^2$ ($=A^2/2$). Verify you get this with the $1/U$ window
correction in place, and show that omitting the correction makes the peak read low
by the window's loss factor.
\end{exercise}

\begin{exercise}
\textbf{Wiener--Khinchin by hand.} For white noise with $r_x[m]=\sigma^2\delta[m]$,
show directly from the Wiener--Khinchin sum that $S_x(\omega)=\sigma^2$ for all
$\omega$ (a flat spectrum). Then for a decaying exponential autocorrelation
$r_x[m]=\sigma^2 a^{\abs{m}}$ with $0<a<1$, compute $S_x(\omega)$ and sketch its
shape. What kind of signal has this ``red'' (low-frequency-heavy) spectrum?
\end{exercise}

\begin{exercise}
\textbf{Overlap payoff.} Estimate the PSD of a fixed record with Bartlett
(non-overlapping) and with Welch ($50\%$ overlap), using the same $L$ and a Hann
window. Welch produces more segments from the same data --- count them. Does the
estimate variance fall by the full ratio of segment counts? Explain why overlapping
segments are not fully independent and what that does to the realized variance
reduction.
\end{exercise}

\begin{exercise}
\textbf{IMU band powers.} Record (or simulate) $20$ s of $100$ Hz accelerometer
data containing a $2$ Hz arm-swing and a $6$ Hz tremor. Compute a Welch PSD with
$L=128$, then integrate it over the bands $0$--$3$ Hz and $4$--$12$ Hz to get two
band-power features. Show how those two numbers change between ``arm swinging'' and
``arm still but tremoring'', and argue why these features are more robust for a
gesture classifier than the raw samples.
\end{exercise}

\begin{exercise}
\textbf{The wrong tool.} You analyze a guitar note's full $4$ s with a single
Welch PSD and get a clean harmonic series, but your pitch tracker still fails during
the note's attack. Explain what Welch averaged away, and which tool from an earlier
chapter you should use instead when the spectrum changes over time.
\end{exercise}

\begin{exercise}
\textbf{Truthful limits.} A teammate proposes selling the wrist controller as a
device that ``reads individual finger movements from the IMU spectrum''. Using the
physics of where finger motion shows up and what a wrist sensor can observe, write
a short, specific rebuttal, and propose what sensing modality would actually be
required for continuous finger pose.
\end{exercise}
